%%=============================================================================
%% Samenvatting
%%=============================================================================

%%---------- Nederlandse samenvatting -----------------------------------------
%
% Deze blok wordt enkel ingevoegd wanneer het document in het Engels wordt
% gecompileerd. Wie de bachelorproef in het Nederlands schrijft, kan dit
% negeren. De inhoud verschijnt dan niet in het document.

\IfLanguageName{english}{%
\selectlanguage{dutch}
\chapter*{Samenvatting}
\lipsum[1-4]
\selectlanguage{english}
}{}

%%---------- Samenvatting -----------------------------------------------------
% De samenvatting in de hoofdtaal van het document

\chapter*{\IfLanguageName{dutch}{Samenvatting}{Abstract}}

Wereldwijd draaien banken, overheden en zorginstellingen op honderden miljarden regels COBOL-code. Tegelijk gaan ervaren COBOL-specialisten met pensioen, terwijl er weinig nieuw talent bijkomt. Dit fenomeen wordt de \textit{silver tsunami} genoemd. Junior developers groeien op met moderne talen zoals Python en met interactieve sandboxes. Voor hen voelt COBOL dubbel vreemd. De syntax is strikt en verbose, en het bestaande leermateriaal is vooral statisch, met PDF-handleidingen en theoretische tutorials. Deze bachelorproef gaat na of \textit{gamification}, het inzetten van game-elementen in een niet-spelcontext, die instapdrempel kan verlagen.

De centrale onderzoeksvraag luidt: \textit{\enquote{Hoe effectief is gamification bij het aanleren van COBOL aan junior developers die gewend zijn aan moderne talen zoals Python?}} Vijf deelvragen behandelen de verschillen tussen COBOL en Python, de leerbarrières voor junioren, de relevante gamification-elementen, het ontwerp van een gamified leeromgeving en de verhouding tot traditionele methoden zoals een papieren cursus.

Het onderzoek verloopt in vier fasen. De eerste fase is een literatuurstudie over de mainframe-context, de verschillen tussen COBOL en Python en de gamification-frameworks PBL, MDA en Octalysis. In de tweede fase wordt een Proof-of-Concept (PoC) gebouwd, \textit{Learn COBOL}. Dat is een interactieve webapplicatie die COBOL probeert aan te leren aan junior developers door gebruik te maken van de geïdentificeerde gamification-elementen in de literatuurstudie. In de derde fase wordt de PoC geëvalueerd in twee delen. Het eerste deel is een vergelijkende literatuurevaluatie, waarin elke ontwerpkeuze wordt afgezet tegen empirisch onderzoek over gamification, programmeeronderwijs en didactiek. Het tweede deel is een uitgewerkt maar niet uitgevoerd pilotontwerp, waarin twee groepen deelnemers elk één dag COBOL leren (sectie~\ref{sec:evaluatie-pilotstudie}), de ene groep via de PoC en de andere groep via traditionele methoden. De vierde fase formuleert de conclusie en de aanbevelingen.

De literatuurevaluatie toont dat de ontwerpkeuzes aansluiten bij wat in de geraadpleegde bronnen staat. De meta-analyse van \textcite{Zhan2022} rapporteert voor gamification in programmeeronderwijs medium-tot-grote effectgroottes (\(SMD \approx 0{,}77\) op motivatie en \(\approx 0{,}75\) op academische prestaties). Omdat de PoC meerdere van die kerncomponenten implementeert (punten, badges, levels, directe feedback en sociale vergelijking), is het aannemelijk dat een pilot effecten in dezelfde grootteorde kan tonen. Dit blijft een verwachting, want er zijn geen metingen op deelnemers uitgevoerd. De PBL-triade wordt verspreid toegepast, in lijn met de waarschuwing van \textcite{KevinWerbach2020} dat leaderboards niet voor iedereen even motiverend werken. In het Octalysis-model van \textcite{Chou2015} domineren de \textit{White Hat}-drijfveren (CD1 tot CD3), die op langetermijnbetrokkenheid mikken. De \textit{Black Hat}-elementen (CD6 tot CD8) worden bewust spaarzaam ingezet.

De toegevoegde waarde voor het werkveld is viervoudig. Er is een theoretisch kader dat gamification (PBL, MDA en Octalysis) koppelt aan mainframe-onderwijs. Er is een functionele Proof-of-Concept op moderne webtechnologie (React, TypeScript, Supabase en Monaco). Er is een uitgevoerde literatuurevaluatie die per ontwerpkeuze onderbouwt waarom het ontwerp aannemelijk is. Tot slot is er een uitgewerkt pilotontwerp dat als blauwdruk voor empirische validatie kan dienen. De belangrijkste beperkingen worden in de bachelorproef expliciet benoemd. Een literatuurevaluatie stelt geen empirische geldigheid voor de PoC zelf vast. De meta-analyses gaan vooral over moderne talen en niet over COBOL. De voornaamste aanbeveling is de uitvoering van het pilotontwerp uit sectie~\ref{sec:evaluatie-pilotstudie}.
